\section{K-022: cancellation between smooth fixed curves}

\begin{theorem}[K-022]
For every $M=2^a\ge8$, there is a positive grading on the fixed
seven-dimensional ring
\[
 R=\mathbb C[\operatorname{Mat}_{2\times5}]^{SL_2}
       =\mathbb C[\operatorname{aGr}(2,5)]
\]
with the following properties.  The ring is a factorial, normal,
Cohen--Macaulay Gorenstein domain with rational singularities,
and $\operatorname{Spec}R$ is smooth away from its vertex.
The nonzero fixed locus of the grading subgroup $\mu_M$ consists
of two smooth density-one curves.  Their local Hilbert residues
are nonzero and opposite, and the Hilbert series of $R$ is regular
at every primitive $M$th root of unity.

Thus generic smoothness of all the relevant fixed components,
even together with these ring properties and an invariant-ring
realization for a connected semisimple group, does not suffice
to transfer a universal-denominator factor to an actual Hilbert
pole.  This is a counterexample to that proposed general transfer,
not a claim of cancellation for the triple-tensor invariant ring.
\end{theorem}

\begin{proof}
Put $m=M/4$ and give the five columns of the matrix the weights
\[
             a=(1,m,2m-1,2m+1,4m-1).                    \tag{K-022a}
\]
They are positive integers.  The Pl\"ucker coordinate $p_{ij}$
has weight $a_i+a_j$.  The invariant-ring identification and
Pfaffian presentation are the classical affine Grassmannian
construction, explicitly stated by Corti--Reid,
\emph{Weighted Grassmannians}, arXiv:math/0206011v1, \S2,
items(iii)--(v), printed p.4.  Their Proposition2.1, equation(2.4),
printed pp.5--6, gives its equivariant free resolution.  Restricting
that resolution to the indicated one-parameter subgroup gives
\[
 H_R(t)=
 \frac{1-\sum_{i=1}^5t^{s-a_i}+\sum_{i=1}^5t^{s+a_i}-t^{2s}}
      {\prod_{1\le i<j\le5}(1-t^{a_i+a_j})},
 \qquad s=\sum_i a_i=9m.                                \tag{K-022b}
\]
This is also the Hilbert-numerator formula in the same source,
printed p.8, with its parameter $u=0$.

We record the geometric properties to exclude a singular-locus
explanation for the cancellation.  On the open set $p_{ij}\ne0$,
the Pl\"ucker equations solve for all remaining coordinates in
terms of $p_{ij}^{\pm1}$ and the six coordinates $p_{ik},p_{jk}$,
$k\notin\{i,j\}$.  This is a Laurent polynomial ring in one
variable and a polynomial ring in six variables.  These charts
cover the punctured cone, which is therefore smooth.

The invariant ring is factorial by the following short argument.
Factor an invariant into irreducibles in the ambient polynomial
ring.  The connected group $SL_2$ preserves every one of the
finitely many irreducible-factor lines.  Each factor is a
semi-invariant, and the corresponding character of $SL_2$ is
trivial.  Consequently every irreducible factor is itself
invariant.  Existence and uniqueness of factorization descend
from the polynomial ring.  In particular $R$ is normal.
The length-three self-dual Pfaffian resolution in Proposition2.1
has a rank-one last term and the ideal has codimension three;
it proves that $R$ is Cohen--Macaulay and Gorenstein.

There is also a direct rational-singularity check.  Use the
ordinary Pl\"ucker grading for this ungraded assertion.  The
resolution on $\mathbb P^9$ is
\[
 0\longrightarrow\mathcal O(-5)\longrightarrow
 \mathcal O(-3)^5\longrightarrow\mathcal O(-2)^5
 \longrightarrow\mathcal O\longrightarrow
 \mathcal O_{\operatorname{Gr}(2,5)}\longrightarrow0.
\]
After twisting by any integer $d\ge0$, every displayed ambient
line bundle has zero higher cohomology: intermediate cohomology
on $\mathbb P^9$ vanishes, and its top cohomology vanishes for
degree at least $-9$.  Breaking the resolution into short exact
sequences yields
$H^q(\operatorname{Gr}(2,5),\mathcal O(d))=0$ for $q>0$.
The total space of $\mathcal O_{\operatorname{Gr}(2,5)}(-1)$ is a
smooth resolution of the affine cone.  Its projection to the
Grassmannian is affine with structure-sheaf pushforward
$\bigoplus_{d\ge0}\mathcal O(d)$, so its higher structure-sheaf
cohomology vanishes.  Since the cone is affine, the higher direct
images for the resolution vanish; normality identifies the
degree-zero direct image with $\mathcal O$.  This proves rational
singularities.  The projective-space cohomology used here is the
standard formula recorded in the Stacks Project, tag01XT.

Exactly two of the ten Pl\"ucker weights are divisible by $M$:
\[
                       a_1+a_5=a_3+a_4=M.
\]
Indeed all weights lie strictly between zero and $2M$, and the
other pair sums from (K-022a) differ from $M$ for $m\ge2$.
On their two-dimensional coordinate plane the Pl\"ucker ideal
reduces to $(p_{15}p_{34})$.  Hence the entire fixed scheme is
\[
               \operatorname{Spec}\mathbb C[p_{15},p_{34}]
                              /(p_{15}p_{34}),           \tag{K-022c}
\]
with both variables of degree $M$.  Its two nonzero curves are
smooth points of the ambient cone by the chart calculation.
Each closure is a polynomial line of weight $M$, so has density
one when the leading term is normalized by $1-t^M$.

Let $\zeta=\exp(2\pi i/M)$ and put $x=\zeta$.
At the line indexed by the pair $\{i,j\}$, the six moving
cotangent eigenvalues, in the point-pullback convention, can be
written as
\[
               \zeta^{a_k-a_i},\quad\zeta^{a_k-a_j},
                           \qquad k\notin\{i,j\}.
\]
The chart weights are equivalently $a_i+a_k,a_j+a_k$, since
$a_i+a_j=M$.  The five factor eigenvalues in our case are
\[
                    x,\quad i,\quad -x^{-1},\quad -x,
                           \quad x^{-1}.
\]
The two moving determinants are therefore
\[
\begin{aligned}
 \det(1-\zeta\mid N^*_{15})
   &=4(1+x^2)(1+x^{-2})(1-ix)(1-ix^{-1})
     =-4i(x+x^{-1})^3,\\
 \det(1-\zeta\mid N^*_{34})
   &=4(1+x^2)(1+x^{-2})(1+ix)(1+ix^{-1})
     = 4i(x+x^{-1})^3.
\end{aligned}                                                   \tag{K-022d}
\]
Neither is zero, because $M\ge8$.  The two local coefficients
from P-F23 are thus
\[
             \frac{i}{4(x+x^{-1})^3},\qquad
            -\frac{i}{4(x+x^{-1})^3}.                    \tag{K-022e}
\]
These are all the positive-dimensional fixed components, and
their sum is zero.

For an independent global verification of regularity, let $P_M$
be the numerator of (K-022b) and write $u=t^m$.  Direct expansion
gives the exact Laurent-polynomial factorization
\[
\begin{split}
 P_M(t)={}&\frac{(u^2-1)(u^2+1)^2}{t}\,
 \bigl(t^2u^5+u^7\\
 &\hspace{18mm}-t(u^{12}-u^{10}+2u^8-2u^6+2u^4-u^2+1)\bigr).
\end{split}                                                   \tag{K-022f}
\]
Since $u^2+1=t^{M/2}+1=\Phi_M(t)$, this supplies two copies of
$\Phi_M$ in the numerator.  The denominator contains exactly two,
coming from the two weights equal to $M$.  Therefore the global
Hilbert series is regular at every primitive $M$th root, in
agreement with the exact cancellation of (K-022e).
\end{proof}

\paragraph{The first example, with its fully factored series.}
For $M=8$ the column weights are $(1,2,3,5,7)$, and the residues
are $i\sqrt2/16$ and $-i\sqrt2/16$.  Define
\[
\begin{split}
 Q(t)={}&1+t+t^2+2t^3+t^4+t^5+2t^6+t^7+3t^8+4t^9+3t^{10}\\
        &+4t^{11}+3t^{12}+t^{13}+2t^{14}+t^{15}+t^{16}
                       +2t^{17}+t^{18}+t^{19}+t^{20}.
\end{split}
\]
Its reduced Hilbert series is
\[
 H_R(t)=\frac{Q(t)}
 {(1-t)^7\Phi_2(t)^5\Phi_3(t)^3\Phi_4(t)^3\Phi_5(t)^2
   \Phi_6(t)^2\Phi_7(t)\Phi_9(t)\Phi_{10}(t)\Phi_{12}(t)}.
                                                               \tag{K-022g}
\]
The exact verifier checks the Pfaffian numerator factorization,
all denominator exponents, the two complete fixed components,
and both moving products in $\mathbb Q[t]/(t^4+1)$.

The grading here weights different matrix columns differently;
it is not the uniform coefficient grading of the triple-tensor
problem.  Moreover other actual pole orders remain large as $M$
grows.  Accordingly K-022 refutes the qualitative smoothness
transfer, and does not refute every possible quantitative relation
between the largest universal and actual denominator orders.
No historical priority is asserted for this specialized weighted
Grassmannian example.
